\section{Norm saturation in the ten integer equation}\label{sec:lucas}
The ten unrestricted variables are
\[
 (t,g_0,c_0,k_0,l,w,x,y,z,m)\in\Z^{10}.
\]
All other symbols in this section are fixed constants or polynomial
abbreviations, not extra quantified coordinates. The literal division-free
arithmetic graph is \pathf{lucas_h1868/integer10.json}.

\subsection{The fixed source and complete equation}
Put $\kappa=2\,134\,940$ and define
\[
 b=6t+4,\quad g=bg_0,\quad k=b(3k_0+2),\quad C=6c_0-1,
 \quad B_0=2^\kappa b^2.
\]
Let $w_i$ and $q_j$ be the fixed lists in Appendix~\ref{app:kernels},
and $D_0,D_1,D_2$ the three explicitly listed sparse polynomials there.
Define
\begin{gather*}
 \ell_0=\sum_{i=1}^{31}B_0^{w_i},\qquad  v_0=\ell_0+g,\\
 J_0=D_2(B_0)v_0^2+D_1(B_0)v_0+D_0(B_0),\qquad 
 T_0=\sum_{j=1}^{22}B_0^{q_j},\\
 N_0=B_0^{117},\qquad N=B_0^{1986},\\
 M_0=N_0-1-(b-1)\ell_0,\\
 S_c=g+N_0(2J_0+T_0),\qquad T_c=M_0+N_0(B_0-2)T_0-2,\\
 X=N(S_c+T_c+1)-S_c-1,\qquad Y=2^{34\,159\,017}bB_0^{2109},\\
 \mathrm{cap}=bB_0^{116}.
\end{gather*}
The source height is $1868$, certified by all 22 exact extraction identities.
Its actual complete coefficient tuple is the retained source-31 tuple.
Then set
\begin{gather*}
 L=lY,\qquad U=2LX,\qquad V=4gwY,\qquad W=bw,\\
 A=U(V+1),\qquad B_1=2X+1,\qquad S=2A-5,\qquad T=3WC-2(W^2-1),\\
 P_h=(LX)UV,\qquad K=X+1+k(2P_h-2),\qquad \\
 Q=(P_h^2-1)K^2+1,\qquad \\
 \Delta=8\mathrm{cap}^3gK^2
 -g^2\{32(C-KL)^2\mathrm{cap}^3+g^2K^2\}.
\end{gather*}
The new Lucas part is
\begin{gather*}
 \xi=zx,\qquad D=(A^2-4)C^2+4,\\
 E_0=C^2\xi,\qquad E=DE_0,\qquad E_2=E_0E=C^4D\xi^2,\\
 F=4(A^2-4)E_2+1,\qquad G=(A/2-1)(F-1)-1,\\
 \Omega=16E,\qquad H=C+F(B_1-C+\Omega y),\\
 I=(G^2-1)H^2+1.
\end{gather*}
Here $A/2=(LX)(V+1)$ is an integer polynomial abbreviation. No division
operation appears in the literal graph. Finally put
\begin{gather*}
 \rho=T+Sy,\qquad \Sigma=\xi S^5,\\
 \Theta=(\rho^3-\rho S^2-S^3)(\rho^2+23S^2),\\
 r=bQ,\qquad a_0=r+(9\Delta-3)(rz)^2,\qquad a_1=DFI,\\
 Z=\Sigma m-\Theta.
\end{gather*}
The completed ordinary polynomial equation is
\begin{equation}\label{eq:ten}
 \boxed{\bigl(Z^2-\Sigma^2(a_0+a_1)\bigr)^2
              -4\Sigma^4a_0a_1=0.}
\end{equation}
Every exponent is a fixed nonnegative integer.

\subsection{The saturation step and reliability theorem}
The preceding construction used $E^2=C^4D^2\xi^2$ in the definitions of
$F$ and $G$. Only those occurrences are replaced by
$E_2=C^4D\xi^2$. The factor $D$ in $E$ and hence in $\Omega$ remains.
Computing $E_0,E,E_2$ needs the same three multiplications as the previous
$E,E^2$ schedule. This lowers degree without adding a variable or operation. Applied alone to
unchanged height-1873 source data, it gives the control equation in
\pathf{lucas/} with degree $1\,248\,364$ and 481 operations. The
recommended equation additionally uses the certified height-1868 target
layout and an exact shared-power plan, reaching $1\,245\,324$ and 480.
All twelve fixed coefficient values are unchanged. The source layout
changes target positions, not variables, residuals or coefficient ports;
its extraction and carry compatibility proof is in
\pathf{packing/certified/PROOF.md}.

Write $u_0(A)=0$, $u_1(A)=1$,
$u_{n+2}(A)=Au_{n+1}(A)-u_n(A)$, with the usual extension to integer indices.

\begin{theorem}[Saturated Lucas reliability]\label{thm:reliability}
Suppose $A$ is even, $|A|>2$, $B$ is odd with
$1<B<|A|/2-1$, $C$ is odd with $|C|>1$, and $\xi\ne0$.
Use the preceding definitions with $B$ replacing $B_1$.
If $DFI$ is a square, then $C=\pm u_B(A)$.
On the residue slice $3\mid A$, $B\equiv3\pmod4$,
$C\equiv-1\pmod3$, the positive choice $C=u_B(A)$ is forced.
\end{theorem}
\begin{proof}
The quantities $D,E_2,F$ are positive, $|G|>1$, and $I>0$.
Because $D\mid E_2$, one has $F\equiv1$, $G\equiv-1$,
$I\equiv1\pmod D$. Moreover
\[
 2G=(A-2)F-A,\qquad 4I\equiv D\pmod F.
\]
Since $F$ is odd, these congruences make $D,F,I$ pairwise coprime.
A square product therefore makes each a square. Choose in the proof only
$d>0$ with $d^2=D$. The Pell/Lucas classification yields indices $s,\mu,t$
with
\[
 C=u_s(A),\qquad 4C^2d\xi=u_\mu(A),\qquad H=u_t(2G).
\]
Odd $C$ makes $s$ odd. Since $|C|>1$, $|s|\ge3$, and Lucas growth
implies $|C|\ge|A|^2-1$. Thus $|4C^2d\xi|>|u_3(A)|$.
Consequently $n=|\mu|>3$ is even, $F=(v_n(A)/2)^2$, and $H,t$ are odd.
The index-separation lemma modulo $F$ gives $s\equiv\pm t\pmod{2n}$.
The rank-divisibility lemma, using $C^2\mid u_n(A)$, gives $C\mid n$.
Hence $s\equiv\pm t\pmod{2C}$.

For the modified expressions,
\[
 \frac{F-1}{\Omega\xi}=\frac{(A^2-4)C^2}{4},\qquad
 \frac{2G+2}{\Omega\xi}=\frac{(A+2)(A-2)^2C^2}{4},
\]
which are integers. It follows that $F\equiv1$ and $2G\equiv-2\pmod\Omega$.
Therefore $H\equiv B\pmod\Omega$, while
$H=u_t(2G)\equiv u_t(-2)=t\pmod\Omega$, because $t$ is odd.
Since $2C\mid\Omega$, we obtain $s\equiv\pm B\pmod{2C}$.
The inequalities $|s|\le|C|$ and $B<|A|<|C|$ exclude a nonzero multiple
of $2C$ in $s-B$ or $s+B$. Thus $s=\pm B$.
Finally $u_{2j+1}(A)\equiv(-1)^j\pmod{A^2}$ gives
$u_B(A)\equiv-1\pmod3$ on the stated slice, which selects the sign.
\end{proof}
The classical classification, index-separation, and rank-divisibility
interfaces are Sun's Lemmas 2.1, 2.2 and Corollary 2.1 in the modified
certificate paper. The bound on $C$ above is derived after classification;
it is not assumed for arbitrary inputs.

This is a guard-dependent improvement reusable whenever the surrounding
arithmetic bridge proves $|C|>1$. It is not an unconditional replacement
of the standalone modified Lucas certificate. The exceptional-index
countercontrol in the verification section makes that boundary explicit.

\subsection{Full interpolation after saturation}
\begin{theorem}[All residue classes remain available]\label{thm:interpolate}
Let $A>2$ be even, $B>0$ odd, $C=u_B(A)$, and $d=v_B(A)>0$, so $D=d^2$.
For every positive integer $N_*$ there is a positive $\xi$ divisible by
$N_*$ such that every class modulo $\xi$ contains a $y$ for which the
modified $D,F,I$ are squares.
\end{theorem}
\begin{proof}
Choose a positive period $n$ of
$M_A=\left(\begin{smallmatrix}A&-1\\1&0\end{smallmatrix}\right)$ modulo
$4N_*C^2d$. Then $4N_*C^2d\mid u_n(A)$. Put
\[
 \xi=\frac{u_n(A)}{4C^2d},\qquad r_0=\frac{v_n(A)}2.
\]
These are positive integers, $N_*\mid\xi$, $n$ is even, $r_0$ is odd,
and $F=r_0^2$. Now $\Omega=4d\,u_n(A)$. The usual
Cayley--Hamilton calculation supplies period $\lambda=4nr_0$ modulo $F$.
Set $q_*=\lambda/\gcd(\lambda,\Omega)$. We claim
$\gcd(q_*,\xi)=1$.
For $p\mid\xi$, $F\equiv1\pmod\xi$ excludes $p\mid r_0$.
The Lucas valuation bound
$\nu_p(u_n(A))\ge\nu_p(n)$, including $p=2$ for even $A$, yields
\[
 \nu_p(\lambda)=\nu_p(4n)\le\nu_p(4u_n(A))\le\nu_p(\Omega).
\]
Every prime power of $\lambda$ supported on $\xi$ is absorbed by $\Omega$.

Given a class $c_*\pmod\xi$, choose arbitrarily large positive $L_*$
with $q_*\mid L_*$ and $L_*\equiv c_*\pmod\xi$.
Let $j_*=B+\Omega L_*$ and $H_0=u_{j_*}(2G)$.
Since $\lambda\mid\Omega L_*$ and $2G\equiv-A\pmod F$,
$H_0\equiv C\pmod F$. Also $H_0\equiv B\pmod\Omega$.
As $F\equiv1\pmod\Omega$, these congruences make
\[
 y=\frac{H_0-C-(B-C)F}{\Omega F}
\]
an integer and give $H=H_0$. The two quotient identities in the previous
proof hold modulo $\Omega\xi$, giving
$H_0\equiv B+\Omega L_*$ and $C+(B-C)F\equiv B$.
After cancellation, $Fy\equiv L_*\pmod\xi$; hence
$y\equiv c_*\pmod\xi$. The Pell identity makes $I$ a square.
\end{proof}
Simply substituting $d\xi$ into the previous certificate would not prove
this statement: interpolation modulo the new $\xi$ is the substantive
compatibility issue resolved here.

\subsection{Soundness for every signed assignment}
The unchanged source and normalization identities give, before any signs
are known,
\[
 b\equiv Q\equiv1\pmod3,\quad \gcd(b,Q)=1,\quad3\mid A,K,
 \quad a_0\equiv1,\ a_1\equiv0\pmod3.
\]
Thus $a_0\ne a_1$. The odd number $S$ is nonzero. The intersective
polynomial $\pi(v)=(v^3-v-1)(v^2+23)$ has no rational root, so
$\Theta\ne0$. Equation~\eqref{eq:ten} therefore forces $\Sigma\ne0$.
The two-square norm identity implies that $a_0,a_1$ are integer squares
and $\Sigma\mid\Theta$: for $q=Z/\Sigma$ the equation gives
\[
 \left(\frac{q^2+a_0-a_1}{2q}\right)^2=a_0,
\]
with a complementary rational square equal to $a_1$; $q=0$ is excluded by
$a_0\ne a_1$. Rational squares of integers have integral roots, so $q$
is integral. The monic homogeneous quintic's prime-valuation argument
then gives $S\mid T$. In particular $\xi,z,x\ne0$.

Factor $a_0=r\{1+(9\Delta-3)rz^2\}$. The two factors are coprime.
The residue $r\equiv1\pmod3$ excludes a negative signed square; hence
$r$ is a positive square. Applying the same argument to $r=bQ$ makes
$b,Q$ positive squares. If $\Delta\le0$, the other factor would be
negative, so $\Delta>0$. The inherited guard bounds imply
$K\ne0$, $0<g<2\mathrm{cap}$. If $l=0$, then $Q=1-K^2\le0$,
a contradiction. The source bounds now yield $X>3b$, $Y>256$,
$|L|\ge Y$ and $|A|\ge2XY>2B_1+2$.

One new estimate removes all exceptional Lucas coordinates:
\[
 \Delta>0,\ g\ge1\quad\Longrightarrow\quad
 4(C-KL)^2<K^2,
\]
so
\[
 |C|>|K|(|L|-1/2)\ge Y-1/2>1.
\]
Theorem~\ref{thm:reliability} now applies, with
$B_1\equiv3\pmod4$, and gives $C=u_{B_1}(A)$.
The recurrence implies $A-2\mid C-B_1$. Restore in the proof only
$h=(C-B_1)/(A-2)$ and invoke the retained reverse binomial bridge.
It decodes the fixed source into a finite ZFC refutation. The old stronger
$D$ factor and its exceptional-index argument have not been silently reused.

\subsection{Completeness and exact degree}
A refutation gives the same canonical arithmetic witnesses as before,
with $A>2$ even, $C=u_{B_1}(A)\equiv5\pmod6$, $r=bQ>0$ square,
$\Delta>0$ and $S\mid T$. Set $c_0=(C+1)/6$.
The integer $\delta=9\Delta-3$ has 3-adic valuation one, so $\delta r$
is positive nonsquare. A positive Pell solution
$u^2-\delta rz^2=1$ makes $a_0=ru^2$ a square.
Apply Theorem~\ref{thm:interpolate} with $N_*=z$, put $x=\xi/z$, and
use intersectivity of $\pi$ to choose the required $y$ class modulo $\xi$.
Then $\Sigma\mid\Theta$ and $a_1$ is a square. Taking
\[
 m=\Theta/\Sigma+\sqrt{a_0}+\sqrt{a_1}
\]
solves~\eqref{eq:ten}. This is an existential equivalence, not a bijection
between all old and new witness tuples.

For exact degree let $a=\deg A=16\,384$; the unchanged source has
$\deg X=7942$, $\deg Y=4219$, $\deg\mathrm{cap}=233$.
The relevant degrees are
\[
\begin{array}{c|rrrrrrrr}
& D&E_2&F&G&\Omega&H&I&a_1\\\hline
\deg&2a+2&2a+10&4a+10&5a+10&2a+6&6a+17&22a+54&28a+66
\end{array}
\]
and $\deg\Sigma=5a+2$. The guard degree $294\,619$ and center degree
$81\,925$ are below their competitors. The dominant final term is
$\Sigma^4a_1^2$, yielding
\[
 \deg P_{10}=4(5a+2)+2(28a+66)=76a+140=1\,245\,324.
\]
Independent nonzero top-coefficient evaluations modulo three primes
certify attainment on the actual fixed-coefficient slice. The final graph
has 300 multiplications, 155 additions and 25 subtractions: 480 operations.
